\PassOptionsToPackage{table,xcdraw}{xcolor}
\documentclass[10pt, a4paper]{article}

\usepackage[T1]{fontenc}
\usepackage[utf8]{inputenc}
\usepackage{tgpagella}
\usepackage{tgcursor}
\usepackage[spanish, es-noshorthands]{babel}
\usepackage{booktabs}
\usepackage{tabularx}
\usepackage{xcolor}
\usepackage[most]{tcolorbox}
\usepackage[margin=1.5cm]{geometry}
\usepackage{fancyhdr}

\setlength{\headheight}{16pt}
\definecolor{ucbblue}{RGB}{0, 51, 102}
\newcolumntype{P}[1]{>{\raggedright\arraybackslash}p{#1}}
\newcolumntype{Y}{>{\raggedright\arraybackslash}X}

\pagestyle{fancy}
\fancyhf{}
\lhead{\textbf{UCB Sede Tarija} | IMT-221}
\chead{Requisitos de Informe Técnico IEEE (Hito 2 y Hito 3)}
\rhead{Semana 13 - Sesión 2}
\renewcommand{\headrulewidth}{0.4pt}

\begin{document}

\begin{center}
    {\Large\textbf{\textcolor{ucbblue}{Requisitos de Estructura Documental y Trazabilidad (Hitos 2 y 3)}}}[cite: 9]
\end{center}

\begin{tcolorbox}[colback=white, colframe=black, boxsep=1pt, top=2pt, bottom=2pt]
    \footnotesize\raggedright
    \textbf{Comunicado Docente Oficial:} Los informes técnicos tipo IEEE constituyen el entregable documental de los Hitos 2 y 3, conservando sus respectivas verificaciones físicas. Durante el montaje se formularán preguntas breves para identificar dificultades y orientar correcciones; estas \textbf{no} son defensas orales formales ni reemplazan el examen individual EC2.\\[0.1cm]
    \textbf{Fecha de Entrega:} Por confirmar. \quad \textbf{Rúbrica:} Por aprobar. \quad \textbf{Formato IEEE:} Aprobado por docente (28/10); \textit{Validez Institucional: Por verificar (reemplazo de F.I.V.E.)}[cite: 9]
\end{tcolorbox}

\section*{1. Estructura Requerida para Informe Hito 2 (Acondicionamiento In-Amp)[cite: 9]}
El reporte debe permitir reproducir lo declarado y separar los resultados teóricos de las hipótesis físicas. No se exige un documento "completo" si el montaje físico no alcanzó la etapa de integración.
\begin{enumerate}
    \setlength{\itemsep}{2pt}
    \item \textbf{Título, autoría y contribuciones:} Grupo, versión, fecha y declaración de responsabilidades.
    \item \textbf{Resumen:} Problema medible, solución propuesta, estado real. No incluir mejoras no verificadas.
    \item \textbf{Problema y Baseline H1:} Esquema/medida disponible (Twin-T), limitación, métrica $M$ y criterio de éxito.
    \item \textbf{Arquitectura H2:} Diagrama completo, interfaces, referencias de nodo, alimentación y motivación del diseño.
    \item \textbf{Método de Diseño:} Topología real, derivación algebraica de ganancia, selección de resistores (matching), rechazo de modo común predicho/medido y justificación de límites con el datasheet del IC real.
    \item \textbf{Simulación LTSpice:} Modelos, amplitudes/fuentes, mediciones identificadas. No atribuir fidelidad física no comprobada.
    \item \textbf{Montaje y Procedimiento:} Fotos etiquetadas, revisión pre-power, instrumentos, carga y secuencia de pruebas.
    \item \textbf{Resultados Físicos:} Datos crudos con unidades, gráficos y cálculos incrementales. Diferenciar claramente etiquetas: \texttt{PREDICHO}, \texttt{SIMULADO}, \texttt{MEDIDO}, \texttt{PENDIENTE}, \texttt{NO OBTENIDO}.
    \item \textbf{Discusión:} Análisis de diferencias teoría/simulación/banco, saturación, incertidumbre y comparación sistémica contra Hito 1 bajo condiciones comparables (o razón documental de no disponibilidad).
    \item \textbf{Conclusiones y Trabajo Pendiente:} Respaldadas exclusivamente por los registros de banco.
\end{enumerate}

\section*{2. Estructura Separada / Identificable para Informe Hito 3 (Filtro Activo)[cite: 9]}
La evidencia del Hito 3 debe reportarse de manera independiente y distinguible del Lab 6.
\begin{itemize}
    \setlength{\itemsep}{2pt}
    \item \textbf{Diseño de Filtro:} Adaptación del esquema a la topología activa (Ej. Butterworth 4.º orden).
    \item \textbf{Respuesta Frecuencial:} Curva de Bode (simulada y medida), análisis por etapas y en cascada.
    \item \textbf{Condición de Interferencia:} Prueba y justificación del criterio de atenuación ($50\,\text{Hz} < -40\,\text{dB}$) preservando la señal útil del sensor de Hito 1/Hito 2.
    \item \textbf{Carga y Estabilidad:} Análisis de interacción de impedancias con el Hito 2.
\end{itemize}
\textit{Aviso: No confundir la curva de un Sallen-Key de segundo orden (Lab 6) con la demostración de desempeño del filtro integrado de cuarto orden (Hito 3).}

\end{document}
